\chapter{Improving Rational Thinking: Training, Feedback and Transfer}
\label{sec:p19-rationality-15-3}

An instruction to avoid a bias, a correct answer after practice, a retained
skill on a new task, and a beneficial life outcome are different achievements.
De Oliveira and Nisbett's Chapter~15.3 of the 2021 \emph{Handbook of
Rationality} argues that reasoning can improve through instruction in
usable rules, varied examples, practice and feedback. It also asks whether
improved test performance benefits people outside those tests
\cite{DeOliveiraNisbett2021Improve}. The argument is reconstructed here with
its cultural and normative qualifications. Two bounded original evidence
records supplement its review: Morewedge and colleagues' game/video
experiments and Mellers and colleagues' forecasting tournament
\cite{MorewedgeEtAl2015Debiasing,MellersEtAl2014ForecastTraining}.

\section{What counts as improvement?}
\label{sec:p19-r153-standards}

\begin{hbdefinition}{def:p19-r153-standards}{Correspondence, coherence and valued outcomes}
In the chapter's working account, \emph{correspondence} concerns agreement
of a judgment with the world or a specified criterion. \emph{Coherence}
concerns the appropriateness of the reasoning process or relations among
judgments under an identified normative standard. A third question concerns
the usefulness of decisions and their experienced consequences for the
person. These criteria require specification; they are not three names for
a single measured score.

A deductively valid argument cannot have true premises and a false
conclusion under its semantics. It can nevertheless have a false conclusion
when a premise is false. An invalid argument can happen to reach a true
conclusion. Thus observed truth does not identify validity, and validity
alone does not establish the truth of one's premises.
\end{hbdefinition}

The source contrasts Western logical standards with dialectical traditions
that emphasize change, contextual balance and socially useful harmony. It
discusses East Asian and Latin American perspectives without claiming that
every member of a society shares one cognitive style. Its historical
explanation contrasts Greek philosophy's role in debate with East Asian
philosophy's social orientation; this is the chapter's proposed cultural
explanation, not an exhaustive history or a tested training mechanism.
A proposition true
at one time and false at another is not thereby both true and false under
one fixed interpretation. The cultural discussion is an attributed account
of reasoning ideals, not a proof that classical consistency is impossible.
The authors explicitly delimit their training review mainly to Western
research and standards \cite[pp.~801--802]{DeOliveiraNisbett2021Improve}.

The ski-ticket example makes value assumptions visible. A person has paid
\$100 for a Michigan trip and \$50 for a Wisconsin trip, cannot sell or
return either ticket, and can take only one because the dates coincide.
Wisconsin is expected to be more enjoyable. With those expenditures
irrecoverable and otherwise equal future costs, choosing the larger
previous expenditure does not recover any money. The economic benchmark
therefore favors Wisconsin. The chapter preserves the objection that
violating an anti-waste rule may itself spoil the experience. If that cost
changes the ranking of \emph{total} future utility, the original enjoyment
comparison no longer establishes the choice. This is a dispute about the
evaluation model, not evidence that past spending becomes recoverable.
The source's report that many untrained respondents choose Michigan has
no sample or interval attached in this chapter; it is not a prevalence
estimate for all reasoners \cite[p.~802]{DeOliveiraNisbett2021Improve}.

\section{Checks on the chapter's illustrative tasks}
\label{sec:p19-r153-source-checks}

The weather argument has propositions \(P\), raining, and \(Q\), wet
streets, under material-conditional semantics. The four-card illustration
has one letter side and one number side per card. Its rule is that a vowel
on the letter side requires an even number on the other side. The visible
cards are A, B, 4 and 7. These are stipulated task interpretations, not
claims about how every participant understood the instructions.

For the source's hypothetical screening question, let \(C\) be cancer,
\(C^c\) its complementary no-cancer event, and
\(+\) a positive test in the specified target population. The given
probabilities are \(\mathbb P(C)=.01\),
\(\mathbb P(+\mid C)=.90\), and
\(\mathbb P(+\mid C^c)=.09\). They are assumed applicable to the
person's population, and cancer/no cancer exhaust the cases. These
numbers are an instructional scenario, not current screening guidance.

\begin{hbtheorem}{thm:p19-r153-example-checks}{proved=annex:p19.r153.example-checks}{Exact conclusions under the stipulated examples}
The premises \(P\to Q\) and \(\neg P\) do not entail
\(\neg Q\). For the four-card task, A and 7 constitute the minimal
fixed inspection set that determines the rule's status for every possible
completion of the hidden sides; B and 4 are unnecessary. An adaptive
inspection can stop as soon as a violation is observed.

Let \(u_M,u_W\) be the total future utilities of the trips and let
\(K\) be their common irrecoverable past cost. If \(u_W>u_M\),
then \(u_W-K>u_M-K\). The inequality need not survive an additional
Wisconsin-specific future disutility.

The screening probability under the stated inputs is
\[
 \mathbb P(C\mid +)=\frac{.01\cdot .90}
  {.01\cdot .90+.99\cdot .09}=\frac{10}{109}
  \mathrel{\approx}.09174.
\]
In a frequency construction with 10000 people, there are 100 cancer cases,
90 true positives, 9900 people without cancer, and 891 false positives:
90 of 981 positives have cancer. This is about 9.17 percent, not exactly
10 percent.
\end{hbtheorem}

The chapter's ``10 percent'' answer is retained as an approximation; the
exact values above give finer arithmetic without changing its base-rate lesson.
Its review reports that, among 160 gynecologists answering a four-choice
question, 21 percent initially answered correctly and 87 percent did so
after a frequency-conversion session. That is the 2021 authors' report of
Gigerenzer and colleagues' 2007 account. The chapter does not supply
assignment, a concurrent control group, item wording for both tests,
exclusions or interval estimates. The change cannot be treated here as a
controlled 66-percentage-point causal effect, retained skill, or improved
patient outcomes \cite[p.~803]{DeOliveiraNisbett2021Improve}.
Chapter~\ref{sec:p19-natural-frequencies-bayesian-reasoning} supplies the
distinct original frequency-format experiments; their samples and stimuli
must not be reassigned to this gynecologist example.

\section{Diagnosis is not yet a remedy}
\label{sec:p19-r153-diagnosis}

The source surveys reasoning failures involving neglect of an actor's
situation, availability instead of base rates or valid predictors,
gain/loss framing, and difficulty testing conditionals. These are its
selected literature examples, not one experiment or one population-wide
error rate. Its history moves from pessimism about judgment to evidence
that relevant rules sometimes exist but are not recruited by a task.
An unsuccessful response may reflect an absent rule, an unrecognized
application, a misunderstood representation or another task goal.
Those possibilities imply different interventions.

Two source illustrations motivate this change of view. A reported failure
to reproduce a sampling-error finding separated people with and without
statistics coursework: the former performed better. Coursework was not
randomly assigned in that illustration, so selection, prior ability and
other educational differences remain possible explanations. Recasting
the card task as checking whether receipts above \$20 have signatures
made relevant violations more recognizable: inspect large receipts and
unsigned backs. This is a task-representation change, not by itself a
demonstration that a person learned a general logical competence.
The chapter cites Wason (1960) for its card illustration, but that
reference's title concerns a conceptual hypothesis task; this treatment
does not invent an original 1960 four-card experiment
\cite[pp.~802--803, 806]{DeOliveiraNisbett2021Improve}.

The authors retain substantial room for improvement in social, medical,
economic and political judgments, including judgments by experts. Yet
expertise is neither a guarantee of correctness nor a measured absence
of usable knowledge. Their practical distinction is between \emph{teaching
new strategies} and \emph{encouraging appropriate use of existing strategies}.
Changing question structure, arranging group deliberation, and replacing
judgment with a mathematical model are identified alternatives; the
chapter selects training for sustained discussion rather than proving
those alternatives inferior \cite[p.~803]{DeOliveiraNisbett2021Improve}.

\section{Practice, transfer, retention and generalized efficacy}
\label{sec:p19-r153-transfer}

\begin{hbdefinition}{def:p19-r153-transfer}{Specify which achievement was tested}
\emph{Task practice} improves responding on practised items or procedures.
\emph{Near transfer} concerns new items with relatively similar operations,
representations or contexts. \emph{Far transfer} concerns a specified more
distant change of domain, representation, purpose or setting. These are
descriptions of a comparison, not a universal binary classification:
record exactly what was preserved and changed.

\emph{Retained improvement} is assessed after a stated interval, with
intervening practice, feedback and exposure recorded. Delayed performance
on practised items is not automatically delayed transfer.
\emph{Generalized efficacy} concerns a specified broader population or
real-world outcome. It requires more than success on selected test items.
\emph{Normative justification} concerns why the scored criterion is
appropriate under the problem's assumptions and values. No empirical
improvement percentage supplies that justification by itself.
\end{hbdefinition}

These distinctions organize the source's questions about training type,
duration of effects and cross-domain use. They also clarify its concluding
gap. A correlation between competence and a good outcome does not show
that teaching competence causes that outcome. Recognizing the word
``confirmation bias'' does not demonstrate reduced confirmation bias;
reduced biased responses on new scenarios do not establish improved
employment, health or relationships. A repeated test, a new sample with
the same protocol, and an off-task transfer test answer different questions.

\section{The statistical and economic training reviewed by the source}
\label{sec:p19-r153-statistical-training}

The restaurant illustration asks why a traveling saleswoman's second meal
often disappoints after an outstanding first meal. The source contrasts
nonstatistical explanations, such as changing chefs or expectations, with
sampling variation and regression after an extreme observation. It reports
statistical answers about 20 percent of the time among people with one
statistics course, 40 percent among new psychology graduate students, and
80 percent among doctoral scientists; respondents without statistical
background gave nonstatistical answers. Those educational groups are not
a randomized training contrast. The chapter supplies neither their sample
counts nor uncertainty for these approximate percentages
\cite[p.~803]{DeOliveiraNisbett2021Improve}.

A statistical explanation is appropriate only when the first observation
is an unusually favorable sample of variable performance and the later
observation is not subject to the same selection. A changing restaurant
or changed expectations can also explain disappointment. An extreme
observation does not logically guarantee a less extreme successor.
Similarly, the ballet-audition example treats observed moves as a sample
of a dancer's larger repertoire. Representativeness, dependence among
moves, selection and genuine changes in ability matter. The mathematical
law of large numbers is not the assertion that every larger sample is
closer to its population than every smaller sample; its actual hypotheses
are retained in Chapter~\ref{sec:p15-weak-laws-and-triangular-arrays}.

The chapter's account of Fong, Krantz and Nisbett (1986) describes five
randomly assigned conditions. One received an abstract statistical packet
and an urn demonstration. Another worked through everyday applications,
including auditions. A third combined packet and examples. Controls
received either no training or definitions alone. Outcomes were the
frequency and quality of statistical explanations on problems ranging
from lotteries to college choice. The full-training condition invoked
statistical principles in about 64 percent of responses, versus fewer
than half in each control; packet-only and examples-only conditions were
intermediate. That denominator is \emph{responses}, not a percentage
of participants who became statistically competent. The review omits
the original participant counts, scoring reliability and intervals;
the percentages are retained at that review-level standing
\cite[p.~804]{DeOliveiraNisbett2021Improve}.

The complementary educational comparison concerns graduate medicine,
psychology and chemistry. The chapter reports increases in conditional
and statistical/methodological reasoning associated with medicine and
psychology, but not chemistry, and relates this to use of statistics and
methodology in everyday examples. This association motivates a teaching
hypothesis; it does not isolate disciplinary training from admissions,
background, curriculum or maturation.

For Fong and Nisbett (1991), the review describes example-based training
in one domain and testing in that or another domain, immediately or after
two weeks. It reports comparable immediate effects across domains, no
reported later decline in the trained domain, and a later decline in the
untrained domain while performance remained above untrained controls.
For Larrick, Morgan and Nisbett (1990), it reports cost--benefit training
with financial or nonfinancial examples, gains on both kinds of measures,
and gains still observed four to six weeks later. Alleged telephone
surveys about politics or consumer preferences gave nonacademic testing
contexts. These are bounded reports of particular rule/domain/time
comparisons, not proof of transfer to every decision. Their original
sample counts, exclusions and uncertainty are not supplied by the 2021
chapter; this record does not fill them with invented values
\cite[p.~804]{DeOliveiraNisbett2021Improve}.

\section{Feedback-rich learning and the original game/video evidence}
\label{sec:p19-r153-games}

The chapter argues that fast, frequent and sufficiently complete feedback
can reveal a reasoning problem and help learners adjust. Its admissions
example requires following accepted \emph{and rejected} applicants over
many cycles; observing only accepted applicants is selective feedback.
Even complete tracking would not directly reveal how the same person
would have performed if admitted to a different institution. Weather
forecasting and sports illustrate repeated feedback opportunities;
hiring and admissions illustrate delayed or incomplete feedback.
These examples motivate a learning-environment design, not an experiment
showing that feedback has one universal effect
\cite[p.~804]{DeOliveiraNisbett2021Improve}.

\begin{hbdefinition}{def:p19-r153-game-record}{Morewedge and colleagues' two intervention experiments}
Experiment~1 recruited 278 people in a Pittsburgh convenience sample,
132 women, mean age 24.5 years, SD 8.52. Its laboratory analysis had
243 completers, 160 game and 83 video; the eight-week online follow-up
had 196, 130 game and 66 video. Experiment~2 recruited 269 people,
155 women, mean age 27.8, SD 12.01; laboratory totals were 238,
156 game and 82 video, and twelve-week follow-up totals 192,
126 game and 66 video. Most had some college education. Laboratory
and follow-up completion were each paid \$30.

Participants were randomly assigned to an interactive game or an
instructional video, without an untreated or attention-matched neutral
control. The approximate durations were sixty and thirty minutes,
respectively. Pretests, immediate posttests and delayed online tests
used counterbalanced versions of bias measures. Bias commission was
scored from 0 to 100; overall commission averaged three bias scores.
Bias knowledge was measured separately. Sample-size plans were declared
to the sponsor, according to the article; a public preregistration is
not established by that statement.
\end{hbdefinition}

Experiment~1 targeted bias blind spot, confirmation bias and fundamental
attribution error. Experiment~2 targeted anchoring, representativeness
and social projection. Bias blind spot compares oneself favorably with
others; confirmation bias favors confirming evidence; fundamental
attribution error underweights situation; anchoring overweights initially
presented information; representativeness uses similarity in place of
appropriate probability information; projection assigns one's own views
or feelings to others. These are the experiments' operational targets,
not diagnoses derived from one wrong answer. The videos defined biases,
illustrated them and
suggested corrective strategies. The games elicited judgments during
missing-person/criminal-investigation narratives, gave personalized
after-action feedback, and supplied further practice and strategies.
Experiment~2 added adaptive extra questions after biased responses.
Game variants and, in Experiment~1, full/first-round versions were pooled
after no significant variant differences were detected; absence of
significance does not prove exact equivalence
\cite[online pp.~4--8, note~2]{MorewedgeEtAl2015Debiasing}.

Exclusions in Experiment~1 were early prototypes (20), crashes (3),
server errors (6), and failure to finish within four hours (6), totaling
35; 47 more missed the follow-up window. Experiment~2 lists crashes
(1), experimenter/participant errors (3), and attention failures (27),
totaling 31. Its note lists 45 later omissions, whereas the printed
238-to-192 totals differ by 46. The unresolved one-person discrepancy
must not be hidden by altering a published denominator. Attrition,
technical exclusions and online versus laboratory administration limit
interpretation \cite[online pp.~4--6, 9--10, note~1]{MorewedgeEtAl2015Debiasing}.

In the statistics below, \(F\) and \(t\) denote the source's reported
comparison statistics with displayed degrees of freedom, \(p\) its
reported significance level, and \(d\) its standardized effect size;
they are not new estimates from reanalysis.
For Experiment~1, overall pre/post changes had reported
\(d=1.68\) for game and \(.69\) for video; pre/follow-up changes
had \(d=1.11\) and \(.66\). Pretest-adjusted game/video contrasts
were \(F(1,240)=68.8,p<.001\) immediately and
\(F(1,193)=12.69,p<.001\) at follow-up. In Experiment~2 the
corresponding change sizes were \(1.74,1.75\) immediately and
\(1.16,1.07\) at follow-up. Its adjusted overall game/video
contrasts were \(F(1,235)=13.44,p<.001\) and
\(F(1,189)=3.66,p=.057\). The latter is not a uniformly strong
delayed game advantage \cite[online pp.~6, 8]{MorewedgeEtAl2015Debiasing}.

The original authors report relative reductions in aggregate bias scores
of 46.25/18.60 percent for game/video immediately in Experiment~1,
and 34.76/20.10 percent at follow-up; Experiment~2 gives
31.94/25.70 percent immediately and 23.57/19.20 percent later.
These are relative score reductions, not percentage-point changes in
the proportion of people who reason rationally. Figures supply 95-percent
confidence-interval bars; the text gives no numerical interval endpoints
for those selected relative changes
\cite[online p.~9, figure~2]{MorewedgeEtAl2015Debiasing}.

The game trained three of six measured confirmation-bias facets in
Experiment~1. On the \emph{untrained} facets, reported game pre/post and
pre/follow-up comparisons were \(t(159)=10.05,p<.001,d=.79\)
and \(t(129)=7.42,p<.001,d=.65\). This concerns transfer among
specific facets and test contexts, not unmeasured daily-life outcomes.
There was no detected game advantage for anchoring in Experiment~2,
or for representativeness at its delayed test. Videos produced better
immediate bias-knowledge scores in both experiments despite the game's
better overall commission result. Knowledge, scored behavior and
retention therefore need separate conclusions
\cite[online pp.~6--8, figures~2--3]{MorewedgeEtAl2015Debiasing}.

The randomized contrast supports comparison of the two assigned packages
in the analysed recruited sample, subject to exclusions and missing
outcomes. Their within-group pre/post improvements alone do not isolate
training from retesting or other time effects. Nor does the package
contrast isolate feedback: game and video also differed in duration,
interaction, narrative, practice and content. The 2021 chapter's
feedback/practice interpretation is a plausible mechanism proposal, not
a separately randomized component effect. The two experiments extend
the approach to different bias sets; they are not an independent
replication by a new team. The main reports do not substitute for the
complete supplemental stimuli, scale-validation record or participant data.

\section{The original geopolitical forecasting record}
\label{sec:p19-r153-forecasting}

Mellers and colleagues studied forecasting during September 2011--April
2012 and June 2012--April 2013. Recruitment through societies, research
centres, alumni associations, science blogs and word of mouth required
at least a bachelor's degree and a roughly two-hour test battery.
Participants were mainly US citizens (76 percent) and men (83 percent),
mean age 36; 64 percent had postgraduate experience. This is not a
representative sample of the world population
\cite[online p.~3]{MellersEtAl2014ForecastTraining}.

Year~1 began with 2246 participants, comprising 1593 survey forecasters
and 653 prediction-market traders. The survey experiment randomly
assigned three training conditions crossed with independent forecasting,
crowd-belief access or teams; market outcomes were not reported in this
article. Training was absent, a roughly 45-minute scenario module, or
a comparable probability module. Scenario instruction encouraged
alternative futures and coherent decision trees. Probability instruction
covered reference classes, combining estimates, time extrapolation,
base rates and biases. Teams of up to 25 could exchange information,
forecasts, rationales and critiques and also received collaboration
instruction. The reported first-year attrition was 7 percent.

Year~2 began with 1648 participants, 943 survey forecasters and 705
market traders. It retained assignments in surviving conditions,
reassigned some returnees, and randomized new recruits. Nominal
no-training included some previous scenario trainees; 86 percent were
returnees. Teams were capped at fifteen. Sixty top first-year performers,
five from each initial condition, were deliberately placed in five elite
teams of twelve. Reported attrition was 3 percent. Consequently the
second year is neither an independent cohort replication nor an
equivalent random contrast of tracking versus no tracking
\cite[online p.~3]{MellersEtAl2014ForecastTraining}.

Forecasters could select questions and update until closure. Eighty-five
questions resolved in Year~1 and 114 in Year~2, 199 in total; most
were binary, with other questions multinomial or conditional. Questions
were open a mean 102 days, range 2--418. The primary accuracy measure
averaged squared forecast/outcome differences over days and questions.
The source's binary score sums \emph{both} outcome components and
therefore ranges from zero to two, twice the single-event squared loss
of Chapter~\ref{sec:p19-calibration-extreme-confidence}. A .9/.1
forecast scores .02 if the first outcome occurs and 1.62 otherwise.
For the treatment comparisons, question-level standardization addressed
differences in question selection; it does not make every choice of
questions or missing forecast causally irrelevant. Payments depended
on minimum participation, not accuracy; leaderboards supplied status
and score feedback \cite[online pp.~3--4]{MellersEtAl2014ForecastTraining}.

Year~1 training had a reported main effect
\(F(2,1586)=14.29,p<.001\). Probability training outperformed
scenario training, \(t(1053)=2.23,p=.026\), which outperformed
no training, \(t(1056)=3.25,p<.001\). Year~2's probability/no-training
comparison was \(F(1,882)=19.12,p<.001\). Group-influence main
effects were \(F(2,1586)=60.68,p<.001\) and
\(F(1,882)=62.76,p<.001\). In Year~1, teams outperformed
crowd-belief forecasters, who outperformed independent forecasters.
Figure~1 supplies two-standard-error bars, not numerical confidence
intervals for every contrast. Year~2 also had a training/group interaction,
\(F(1,882)=5.78,p=.018\); one overall training benefit must not
erase that dependence \cite[online pp.~4--5]{MellersEtAl2014ForecastTraining}.

The actual Table~1 gives an illustrative outcome contrast: among
questions open at least a month, Year~1 first-week mean scores were
.44 for untrained independent forecasters and .35 for probability-trained
teams; last-week means were .31 and .19. These are raw cell means on
a selected question subset, not the standardized main-effect estimates
or an isolated effect of feedback. Training and teaming differences
persisted in regressions adjusting for initial forecast timing and
update counts; that adjustment does not identify every mediating process.

The record includes improved calibration and resolution, evaluated
separately from mean accuracy. Its repeated event forecasts are not
independently recruited participants. It also contains a failure case:
on the unexpected fatal maritime-confrontation question, trained teams'
average score was 1.44 versus 1.31 for untrained independent forecasters.
Average advantage does not imply advantage on every event. Superior
second-year elite-team performance supports persistence of selected
forecasters' success, but previous selection, effort and collaboration
prevent an isolated causal attribution to tracking
\cite[online pp.~5--9]{MellersEtAl2014ForecastTraining}.

This tournament provides real-event evidence over extended periods,
with continuing opportunity to practise, revisit modules and receive
feedback. It does not establish retention without reinforcement or
transfer to clinical, financial or personal decisions. Neither the
2021 chapter's visual impression that teaming helped more than training
nor comparing their significance statistics establishes a universal
ranking of those interventions. Their packages, contrasts and costs
differ.

\section{Choosing an intervention and stating its warranted conclusion}
\label{sec:p19-r153-method}

\begin{hbmethod}{def:p19-r153-intervention-method}{A rule-and-transfer intervention record}
Specify the target judgment and justify its answer criterion. Determine
whether the obstacle is missing knowledge, recognition of applicability,
representation, monitoring or control. Teach the rule with explicit
conditions, contrasting applications and failure cases; use new practice
items and feedback that distinguishes incorrect reasoning from an
unlucky outcome. Record who supplies feedback, how quickly and whether
unsuccessful or unselected cases remain observable.

Compare the intervention with a specified alternative under a suitable
assignment design. Match time and attention when isolating a component
such as feedback. Assess bias knowledge, task performance, confidence,
transfer and delayed outcomes separately. Record exclusions, attrition,
task familiarity, continued exposure, scoring, dependence and uncertainty
using \ref{def:p19-evidence-record}. State what a delayed or new-domain
test actually changes. Do not relabel repeated practice as independent
replication, or a tested package as proof of one psychological mechanism.
\end{hbmethod}

The interfaces matter. Chapter~\ref{sec:p19-rationality-2-5} distinguishes
mindware gaps, failed engagement and focal representations: an instruction
to try harder may repair none of them. Chapter~\ref{sec:p19-rationality-8-6}
separates object-level correctness from confidence and control: learning
to recognize a bias need not produce a calibrated feeling of rightness
or a useful stopping decision. Chapter~\ref{sec:p17-metareasoning} supplies
conditional computational-control models, not evidence that one human
training programme optimizes all deliberation costs.

Chapter~\ref{sec:p19-analogy-and-problem-transfer} distinguishes retrieval,
noticing relevance, mapping and adaptation, and shows why a hint matters.
The separately allocated Oxford 2012 treatment \emph{Learning to Think:
Cognitive Mechanisms of Knowledge Transfer} addresses knowledge components,
cognitive task analysis, instructional redesign and symbolic languages.
That is a different explanatory presentation from this chapter's selected
statistical, debiasing and forecasting interventions. A successful worked
analogy does not by itself demonstrate durable rule transfer; neither
presentation discharges the other's source-specific evidence obligations.

For subject/task lookup, a suitable success case is an identified sampling
rule practised across contrasting everyday examples, with new-domain
questions and a delayed control comparison. The reported Fong study is
evidence at its stated review level, not a measured guarantee for a new
classroom. A clear failure case is a programme that merely teaches bias
names, retests the identical quiz immediately, and then claims lasting
general rationality. The game experiments show why naming accuracy and
bias commission must be measured separately. A second failure is
following only accepted applicants and interpreting their success as a
complete evaluation of selection skill.

A concrete training contrast can use the same conditional relation in
the weather, card and receipt contexts. A learner who repeats the rain
answer after memorizing it has demonstrated only that response. A learner
who checks an unsigned receipt for an amount above the threshold can
explain the relevant violation in a new context. Testing that response
without a hint and again after a stated delay distinguishes applicability
and retention from immediate prompted practice. This is an instructional
construction illustrating the comparison, not a reported new experiment
or evidence that the design has already produced transfer.

\begin{proofboundary}{def:p19-r153-evidence-boundary}{Measured improvement is not an unrestricted prescription}
The source's conclusion favors applicable rules, varied examples,
practice and feedback over bare exhortation. Its reported statistical
and cost--benefit studies concern particular domains and delays. The
original game/video and forecasting records strengthen specific evidence
claims while retaining their controls, selection and measurement limits.
No universal effect size, training duration or population benefit follows.

The source finally cites associations of normative competence with less
risky behavior and delinquency, higher salaries among professors and
higher student grades. Those are not randomized effects of this chapter's
training programmes. Baseline ability, opportunity, selection and other
causes remain possible. The authors' summary calls for causal research
on outcomes over time. This remains an open evidential step, not an
established benefit hidden behind the phrase ``good thinking''
\cite[pp.~801, 805]{DeOliveiraNisbett2021Improve}.
\end{proofboundary}
